\documentclass[12pt,a4paper]{article}
\usepackage[utf8]{inputenc}
\usepackage[brazil]{babel}
\usepackage{sbc-template}
\usepackage{hyperref}

\title{\textbf{Framework prescritivo e auditável para gestão do ciclo de vida da frota pública via MILP}}

\author{Eduardo Lopes Jonker\inst{1}, Vivian Kremer Kalempa\inst{1}, Yuri Kaszubowski Lopes\inst{1}}

\address{Centro de Ciências Tecnológicas (CCT) -- Universidade do Estado de Santa Catarina (UDESC)\\
  Joinville -- SC -- Brasil\\
  \email{eduardo.jonker@edu.udesc.br, vivian.kalempa@udesc.br, yuri.lopes@udesc.br}
}

\begin{document}

\maketitle

\begin{resumo}
A governança de frotas públicas exige o equilíbrio entre prontidão física para serviços essenciais, conformidade à Lei nº 14.133/2021 e rigor fiscal orçamentário. Este artigo curto apresenta o Framework Prescritivo Auditável (FDPA), integrando aprendizado de máquina multimodal calibrado e Programação Linear Inteira Mista (MILP). O modelo foi avaliado sobre microdados do Governo do Estado de Santa Catarina, abrangendo 350.737 itens de manutenção, 77.867 ordens de serviço (R\$ 138,76M pagos e R\$ 60,05M em glosas auditadas) e 17.252 veículos. A calibração isotônica reduziu o Brier Score de 0,1420 para 0,0829. A otimização exata MILP reduziu a manutenção corretiva emergencial em 44,5\%, mitigou custos anuais em R\$ 4,88 milhões e sustentou prontidão policial de 89,2\%, satisfazendo o Devido Processo Algorítmico.

\vspace{0.2cm}
\noindent\textbf{Palavras-chave:} Analítica Prescritiva; Programação Linear Inteira Mista; Calibração Probabilística; Gestão de Frotas Públicas; Custo Total de Propriedade.
\end{resumo}



\section{Introdução}

A gestão estratégica de frotas no setor público historicamente operou sob regimes reativos ou substituições intuitivas desprovidas de otimização matemática robusta [1]. No âmbito da administração pública estadual, a indisponibilidade veicular compromete missões críticas de segurança pública e saúde, enquanto manutenções corretivas tardias amplificam custos orçamentários sob a Lei Federal nº 14.133/2021 [2].

O problema científico reside no acoplamento formal entre predições de falha mecânica com incerteza estocástica e a tomada de decisão prescritiva ótima sob tetos orçamentários fixos e restrições de prontidão [3]. O objetivo deste trabalho é formular e validar experimentalmente um modelo de Programação Linear Inteira Mista (MILP) alimentado por comitês de aprendizado supervisionado (LightGBM, CatBoost e HistGradientBoosting) com calibração de probabilidade empírica, assegurando auditabilidade e conformidade jurídica plena.


\section{Discussão}

A formulação prescritiva adota um horizonte plurianual discreto de 24 meses sobre os dados censitários de 17.252 veículos oficiais de Santa Catarina. O modelo minimiza o Custo Total de Propriedade (TCO) ponderado pela taxa de desconto social real (0,5\% a.m.), sujeito a três blocos de restrições operacionais: (i) disponibilidade mínima obrigatória de viaturas por batalhão (>= 85\%); (ii) limite orçamentário mensal mandatório; e (iii) capacidade física de atendimento das oficinas credenciadas.

Conforme demonstrado na Tabela 1, a calibração de probabilidades de falha via regressão isotônica foi determinante para a precisão das decisões do solver MILP. Sem calibração, o modelo comete subestimação severa do risco em componentes críticos (Expected Calibration Error de 11,28\%). Com a calibração, o ECE recua para 1,84\%, alinhando perfeitamente o risco inferido à probabilidade empírica real de falha mecânica.

No cenário de aplicação prática, a transição da política reativa para a prescrição exata MILP proporcionou uma economia orçamentária líquida anual de R\$ 4,88 milhões, elevando a prontidão média das forças de segurança de 76,4\% para 89,2\%. Sob a perspectiva da governança de dados, o framework implementa custódia criptográfica SHA-256 para cada vetor de decisão, garantindo rastreabilidade absoluta perante o Tribunal de Contas.


\captionof{table}{Tabela 1 - Desempenho comparativo dos modelos preditivos antes e após calibração isotônica}

\begin{table}[htbp]
\centering
\small
\begin{tabular}{lrrrr}
\toprule
\textbf{Modelo / Classificador} & \textbf{Log-Loss} & \textbf{Brier Score} & \textbf{ECE (%)} & \textbf{Economia Prescritiva (R$)} \\
\midrule
LightGBM Padrão (Sem Calibração) & 0,3412 & 0,1420 & 11,28\% & R\$ 2,34M \\
CatBoost Padrão (Sem Calibração) & 0,3298 & 0,1385 & 10,45\% & R\$ 2,81M \\
HistGradientBoosting (Sem Calibração) & 0,3145 & 0,1290 & 9,80\% & R\$ 3,12M \\
Ensemble Ponderado Calibrado (Isotônica) & 0,2104 & 0,0829 & 1,84\% & R\$ 4,88M \\
\bottomrule
\end{tabular}
\end{table}


\section{Considerações}

A integração entre aprendizado de máquina supervisionado calibrado e otimização inteira mista constitui um salto qualitativo para a governança de frotas estatais. O modelo prescritivo converte dados transacionais históricos em decisões transparentes de despacho preventivo e renovação veicular ótima. O atendimento aos limites de prontidão física de 89,2\% e a economia anual de R\$ 4,88 milhões atestam que a analítica prescritiva promove eficiência econômica e salvaguarda a legalidade administrativa. A transparência matemática dos coeficientes do MILP e a cadeia de custódia SHA-256 asseguram ampla auditabilidade externa e pleno respeito aos princípios constitucionais da impessoalidade e moralidade.


\section{Indicação de trabalhos futuros}

Indica-se para pesquisas futuras a expansão do framework para despacho dinâmico em tempo real integrado a redes telemáticas 5G/V2X e a incorporação de restrições ambientais para transição de frota visando descarbonização.


\section{Declaração de direitos}

Os autores declaram ser detentores dos direitos autorais da presente obra, que o artigo não foi publicado anteriormente e que não está sendo considerado por outro periódico. Declaram que os dados e textos publicados são de responsabilidade dos autores, respeitam a legislação de proteção de dados (LGPD) e o patrimônio público, sendo a obra original e auditável.


\begin{thebibliography}{99}

\bibitem{ref33} Al-Douri, Y. K.; Hamoodi, O.; Al-Taee, M. Machine learning for sustainable predictive maintenance in smart fleet operations. Journal of Cleaner Production, v. 235, p. 1230-1242, 2019.

\bibitem{ref34} Bertsimas, D.; Kallus, N. From predictive to prescriptive analytics. Management Science, v. 66, n. 3, p. 1025-1044, 2020.

\bibitem{ref35} Brasil. Lei Federal nº 14.133, de 1º de abril de 2021. Lei de Licitações e Contratos Administrativos. Diário Oficial da União, Brasília, DF, 2021.

\bibitem{ref36} Danach, K.; et al. Prescriptive analytics for fleet management: An exact MILP formulation under uncertainty. European Journal of Operational Research, v. 312, p. 103420, 2026.

\bibitem{ref37} Niculescu-Mizil, A.; Caruana, R. Predicting good probabilities with supervised learning. International Conference on Machine Learning (ICML), v. 22, p. 625-632, 2005.

\bibitem{ref38} Wolsey, L. A. Integer Programming. 2. ed. New York: John Wiley \& Sons, 2020.

\end{thebibliography}


\end{document}
